\begin{apartats}

\begin{tria}{extrems-tasca}
\itemtria{original}{1}{1,25}
La funció
\[
f(x)=ax^3+bx^2
\]
té un extrem relatiu en el punt $(2,4)$. Troba $a$ i $b$, digues de quin tipus d'extrem es
tracta i troba l'altre extrem relatiu de $f$.

\begin{solucio}
\textbf{Passa per $(2,4)$}: $8a+4b=4$. \quad \textbf{Extrem en $x=2$}: $f'(x)=3ax^2+2bx$ i
$f'(2)=12a+4b=0$.\\
Restant les dues equacions: $4a=-4$, d'on $\boxed{a=-1}$, i llavors $-8+4b=4$, és a dir
$\boxed{b=3}$.\\
Amb $f(x)=-x^3+3x^2$: $f'(x)=-3x^2+6x=-3x(x-2)$, que s'anul·la en $x=0$ i $x=2$.\\
$f$ decreix a $(-\infty,0)$, creix a $(0,2)$ i decreix a $(2,+\infty)$: en $(2,4)$ hi ha un
\textbf{màxim relatiu} i en $(0,0)$, un \textbf{mínim relatiu}.
\end{solucio}

\itemtria{beneficis-context}{1}{1,25}
Els beneficis d'una empresa, en milers d'euros, $t$ anys després d'obrir, són
\[
B(t)=at^3+bt^2+ct .
\]
El primer any els beneficis van ser màxims, de $8$ milers d'euros, i el segon any la corba va
tenir un punt d'inflexió. Troba $a$, $b$ i $c$. Quan tornen a ser nuls els beneficis?

\begin{solucio}
$B'(t)=3at^2+2bt+c$ i $B''(t)=6at+2b$. Les tres dades són tres equacions: inflexió en $t=2$,
$12a+2b=0$; màxim en $t=1$, $3a+2b+c=0$; i valor màxim, $a+b+c=8$.\\
De la primera, $b=-6a$; de la segona, $c=9a$; i la tercera dona $4a=8$. Per tant,
$\boxed{a=2}$, $\boxed{b=-12}$ i $\boxed{c=18}$: $B(t)=2t^3-12t^2+18t=2t(t-3)^2$.\\
Comprovació: $B''(1)=-12<0$, i en $t=1$ hi ha un màxim. Els beneficis tornen a ser nuls al cap
de $\boxed{t=3}$ anys.
\end{solucio}
\end{tria}

\begin{nomesllarg}
\apartat{0,75}
\textbf{Inventa.} Escriu una funció polinòmica que tingui exactament dos extrems relatius: un
màxim en $x=0$ i un mínim en $x=4$. Justifica-ho.

\begin{solucio}
Cal que $f'$ s'anul·li en $0$ i en $4$, i que hi canviï de positiva a negativa i de negativa a
positiva, respectivament. La paràbola $f'(x)=3x(x-4)$ ho compleix: és positiva abans
de $0$, negativa entre $0$ i $4$ i positiva després.\\
Integrant: $f(x)=x^3-6x^2$.\\
Comprovació: $f'(x)=3x^2-12x=3x(x-4)$, amb màxim en $x=0$ i mínim en $x=4$.
\end{solucio}
\end{nomesllarg}

\apartat[1,25]{0,75}
Pot una funció polinòmica de grau $3$ tenir exactament \textbf{un} extrem relatiu? Raona-ho.

\begin{solucio}
Si $f$ té grau $3$, $f'$ és de grau $2$ i pot tenir dues arrels diferents, una de doble o cap.\\
Amb dues arrels diferents, $f'$ canvia de signe dues vegades: \textbf{dos extrems}.\\
Amb una arrel doble, $f'$ no canvia de signe: \textbf{cap extrem} (per exemple, $f(x)=x^3$).\\
Sense arrels reals, $f'$ tampoc no canvia de signe: \textbf{cap extrem}.\\
Per tant, \textbf{no és possible} tenir-ne exactament un.
\end{solucio}

\end{apartats}
